% ═══════════════════════════════════════════════════════════════
% LERNBLATT: El tiempo y los lugares
% Spanisch Klasse 8 - Unidad 6, DS 1
% ═══════════════════════════════════════════════════════════════

\documentclass[11pt, a4paper]{article}

\usepackage[utf8]{inputenc}
\usepackage[T1]{fontenc}
\usepackage[ngerman]{babel}

% --- SCHRIFTEN ---
\usepackage{helvet}
\renewcommand{\familydefault}{\sfdefault}
\usepackage{palatino}
\newcommand{\seriftext}[1]{{\fontfamily{ppl}\selectfont #1}}

\usepackage{microtype}
\usepackage[margin=1.8cm, top=2cm, bottom=1.5cm]{geometry}
\usepackage{enumitem}
\usepackage{tabularx}
\usepackage{booktabs}
\usepackage{array}
\usepackage{multicol}
\usepackage{tikz}

\usepackage{fancyhdr}
\usepackage{lastpage}

% ═══════════════════════════════════════════════════════════════
% VARIABLEN
% ═══════════════════════════════════════════════════════════════

\newcommand{\lektion}{06a}
\newcommand{\thema}{El tiempo y los lugares}
\newcommand{\untertitel}{Das Wetter und die Orte}
\newcommand{\klasse}{8}
\newcommand{\lernziele}{Wetter beschreiben · Orte im Viertel benennen · Aktivitäten mit Wetter verbinden}

% ═══════════════════════════════════════════════════════════════
% FARBEN
% ═══════════════════════════════════════════════════════════════

\usepackage{xcolor}

\definecolor{colorrojo}{HTML}{C62828}
\definecolor{colorazul}{HTML}{1565C0}
\definecolor{colorverde}{HTML}{2E7D32}
\definecolor{coloramarillo}{HTML}{F9A825}
\definecolor{colorgris}{HTML}{757575}
\definecolor{colornegro}{HTML}{222222}

\definecolor{hellgrau}{HTML}{F5F5F5}
\definecolor{rahmen}{HTML}{CCCCCC}
\definecolor{akzent}{HTML}{666666}

% ═══════════════════════════════════════════════════════════════
% BOXEN
% ═══════════════════════════════════════════════════════════════

\usepackage{tcolorbox}
\tcbuselibrary{skins}

\newtcolorbox{lernzielbox}{
    colback=hellgrau, colframe=hellgrau, boxrule=0pt, arc=0pt,
    left=10pt, right=10pt, top=8pt, bottom=8pt,
    before skip=6pt, after skip=8pt
}

\newtcolorbox{grammatikbox}[1][]{
    colback=white, colframe=white, boxrule=0pt, arc=0pt,
    left=10pt, right=6pt, top=6pt, bottom=6pt,
    borderline west={3pt}{0pt}{akzent},
    before skip=4pt, after skip=4pt
}

\newtcolorbox{merkbox}[1][]{
    colback=hellgrau, colframe=hellgrau, boxrule=0pt, arc=0pt,
    left=10pt, right=10pt, top=8pt, bottom=8pt,
    borderline north={2pt}{0pt}{colornegro},
    before skip=6pt, after skip=6pt
}

\newtcolorbox{vokabelbox}[1][]{
    colback=white, colframe=rahmen, boxrule=0.5pt, arc=0pt,
    left=8pt, right=8pt, top=6pt, bottom=6pt,
    before skip=4pt, after skip=4pt
}

\newtcolorbox{tippbox}{
    colback=white, colframe=white, boxrule=0pt, arc=0pt,
    left=8pt, right=4pt, top=3pt, bottom=3pt,
    borderline west={2pt}{0pt}{akzent}
}

% ═══════════════════════════════════════════════════════════════
% BEFEHLE
% ═══════════════════════════════════════════════════════════════

\newcommand{\es}[1]{\seriftext{\textbf{#1}}}
\newcommand{\de}[1]{\textit{\textcolor{akzent}{#1}}}
\newcommand{\gram}[1]{\textsc{#1}}
\newcommand{\abmark}[1]{{\footnotesize\textcolor{akzent}{$\rightarrow$ AB #1}}}
\newcommand{\abschnitt}[2]{\noindent\textbf{\large #1. #2}}

\setlength{\parindent}{0pt}
\setlength{\parskip}{5pt}

% ═══════════════════════════════════════════════════════════════
% KOPF- UND FUSSZEILE
% ═══════════════════════════════════════════════════════════════

\pagestyle{fancy}
\fancyhf{}
\renewcommand{\headrulewidth}{0.5pt}
\renewcommand{\footrulewidth}{0pt}
\lhead{\footnotesize\textsc{Spanisch} · Klasse \klasse}
\rhead{\footnotesize\thema}
\cfoot{\footnotesize\thepage\,/\,\pageref{LastPage}}

% ═══════════════════════════════════════════════════════════════
% DOKUMENT
% ═══════════════════════════════════════════════════════════════

\begin{document}

% --- TITEL ---
\begin{center}
    {\LARGE\bfseries \thema}\\[3pt]
    {\small \untertitel}
\end{center}

\vspace{4pt}

% --- EINLEITUNG ---
Stell dir vor, du bist in Spanien und planst deinen Tag.
Das Wetter bestimmt, was du machst: Bei Sonne gehst du an den Strand, bei Regen ins Kino.
Heute lernst du, über das Wetter zu sprechen und Orte in deiner Stadt zu benennen.

\vspace{2pt}

% --- LERNZIELE ---
\begin{lernzielbox}
\textbf{Das lernst du:}\quad \lernziele
\end{lernzielbox}

% ═══════════════════════════════════════════════════════════════
% ZWEI SPALTEN
% ═══════════════════════════════════════════════════════════════

\begin{multicols}{2}

% ---------------------------------------------------------------
% ABSCHNITT 1: EL TIEMPO
% ---------------------------------------------------------------
\abschnitt{1}{El tiempo} \de{-- Das Wetter}

\vspace{4pt}

\begin{grammatikbox}
\textbf{Die Frage:}\\[2pt]
\es{¿Qué tiempo hace?} -- \de{Wie ist das Wetter?}
\end{grammatikbox}

\vspace{4pt}

\textbf{Antworten mit \es{hace}:}

\begin{tabular}{@{}ll@{}}
\es{Hace sol.} & \de{Es ist sonnig.} \\[0.4em]
\es{Hace calor.} & \de{Es ist heiß.} \\[0.4em]
\es{Hace frío.} & \de{Es ist kalt.} \\[0.4em]
\es{Hace viento.} & \de{Es ist windig.} \\[0.4em]
\es{Hace buen tiempo.} & \de{Das Wetter ist gut.} \\[0.4em]
\es{Hace mal tiempo.} & \de{Das Wetter ist schlecht.} \\[0.4em]
\es{Hace 20 grados.} & \de{Es sind 20 Grad.} \\
\end{tabular}

\vspace{6pt}

\textbf{Antworten mit Verb:}

\begin{tabular}{@{}ll@{}}
\es{Llueve.} & \de{Es regnet.} \\[0.4em]
\es{Nieva.} & \de{Es schneit.} \\[0.4em]
\es{Hay nubes.} & \de{Es ist bewölkt.} \\
\end{tabular}

\vspace{4pt}

\begin{tippbox}
\footnotesize\textit{Merke:} \es{Hace} + Substantiv (sol, frío, calor...)\\
Aber: \es{Llueve} / \es{Nieva} sind eigenständige Verben!
\end{tippbox}

\abmark{Aufgabe 1}

\columnbreak

% ---------------------------------------------------------------
% ABSCHNITT 2: LOS LUGARES
% ---------------------------------------------------------------
\abschnitt{2}{Los lugares} \de{-- Die Orte}

\vspace{4pt}

\textbf{Orte in der Stadt / im Viertel:}

\begin{tabular}{@{}ll@{}}
\es{el parque} & \de{der Park} \\[0.4em]
\es{la plaza} & \de{der Platz} \\[0.4em]
\es{el cine} & \de{das Kino} \\[0.4em]
\es{el teatro} & \de{das Theater} \\[0.4em]
\es{la biblioteca} & \de{die Bibliothek} \\[0.4em]
\es{la iglesia} & \de{die Kirche} \\[0.4em]
\es{el supermercado} & \de{der Supermarkt} \\[0.4em]
\es{la panadería} & \de{die Bäckerei} \\[0.4em]
\es{la farmacia} & \de{die Apotheke} \\[0.4em]
\es{el restaurante} & \de{das Restaurant} \\[0.4em]
\es{la cafetería} & \de{das Café} \\[0.4em]
\es{el centro comercial} & \de{das Einkaufszentrum} \\
\end{tabular}

\vspace{6pt}

\begin{tippbox}
\footnotesize\textit{Tipp:} Viele Orte auf \es{-ería} sind Geschäfte:\\
\es{panadería} (Bäckerei), \es{carnicería} (Metzgerei), \es{librería} (Buchhandlung)
\end{tippbox}

\abmark{Aufgabe 1}

\end{multicols}

% ═══════════════════════════════════════════════════════════════
% SEITE 2
% ═══════════════════════════════════════════════════════════════
\newpage

% ---------------------------------------------------------------
% ABSCHNITT 3: WETTER + AKTIVITÄT
% ---------------------------------------------------------------
\abschnitt{3}{Cuando hace... voy a...} \de{-- Wenn es ... ist, gehe ich ...}

\vspace{4pt}

\begin{merkbox}
\textbf{Die Struktur}

\vspace{4pt}
\begin{center}
{\large \es{Cuando} + \textit{Wetter}, + \textit{Aktivität}}
\end{center}

\vspace{4pt}
\textit{Beispiel:} \es{Cuando hace sol, voy al parque con mis amigos.}\\
\de{Wenn es sonnig ist, gehe ich mit meinen Freunden in den Park.}
\end{merkbox}

\vspace{4pt}

\begin{multicols}{2}

\textbf{Beispiele:}

\begin{itemize}[leftmargin=1.2em, itemsep=2pt]
    \item \es{Cuando hace calor, voy a la piscina.}\\
          \de{Wenn es heiß ist, gehe ich ins Schwimmbad.}
    \item \es{Cuando llueve, voy al cine.}\\
          \de{Wenn es regnet, gehe ich ins Kino.}
    \item \es{Cuando hace frío, me quedo en casa.}\\
          \de{Wenn es kalt ist, bleibe ich zu Hause.}
    \item \es{Cuando nieva, voy a la montaña.}\\
          \de{Wenn es schneit, gehe ich in die Berge.}
\end{itemize}

\columnbreak

\textbf{Nützliche Verben:}

\begin{tabular}{@{}ll@{}}
\es{voy a...} & \de{ich gehe zu...} \\[0.4em]
\es{voy al parque} & \de{ich gehe in den Park} \\[0.4em]
\es{voy a la playa} & \de{ich gehe an den Strand} \\[0.4em]
\es{me quedo en casa} & \de{ich bleibe zu Hause} \\[0.4em]
\es{juego al fútbol} & \de{ich spiele Fußball} \\[0.4em]
\es{leo un libro} & \de{ich lese ein Buch} \\
\end{tabular}

\end{multicols}

\abmark{Aufgabe 1 + 2}

\vspace{8pt}

% ---------------------------------------------------------------
% ABSCHNITT 4: WETTER + KLEIDUNG
% ---------------------------------------------------------------
\abschnitt{4}{¿Qué ropa llevas?} \de{-- Welche Kleidung trägst du?}

\vspace{4pt}

\begin{grammatikbox}
\textbf{Die Frage:}\\[2pt]
\es{¿Qué ropa llevas cuando hace frío/calor/...?}\\
\de{Welche Kleidung trägst du, wenn es kalt/heiß/... ist?}
\end{grammatikbox}

\vspace{4pt}

\begin{multicols}{2}

\textbf{Beispielantworten:}

\begin{itemize}[leftmargin=1.2em, itemsep=2pt]
    \item \es{Cuando hace frío, llevo un abrigo y una bufanda.}\\
          \de{Wenn es kalt ist, trage ich einen Mantel und einen Schal.}
    \item \es{Cuando hace calor, llevo pantalones cortos y una camiseta.}\\
          \de{Wenn es heiß ist, trage ich kurze Hosen und ein T-Shirt.}
    \item \es{Cuando llueve, llevo botas y un paraguas.}\\
          \de{Wenn es regnet, trage ich Stiefel und einen Regenschirm.}
\end{itemize}

\columnbreak

\textbf{Kleidung (Wiederholung):}

\begin{tabular}{@{}ll@{}}
\es{el abrigo} & \de{der Mantel} \\[0.3em]
\es{la bufanda} & \de{der Schal} \\[0.3em]
\es{los guantes} & \de{die Handschuhe} \\[0.3em]
\es{el gorro} & \de{die Mütze} \\[0.3em]
\es{las botas} & \de{die Stiefel} \\[0.3em]
\es{el paraguas} & \de{der Regenschirm} \\[0.3em]
\es{las gafas de sol} & \de{die Sonnenbrille} \\
\end{tabular}

\end{multicols}

\abmark{Aufgabe 2}

\vspace{8pt}

% ═══════════════════════════════════════════════════════════════
% KULTURELLES
% ═══════════════════════════════════════════════════════════════

\begin{grammatikbox}
\textbf{¿Sabías que...?} \de{-- Wusstest du?}

\vspace{4pt}
\begin{multicols}{2}
\footnotesize

$\bullet$ In Spanien sagt man oft \es{hace un frío que pela} (es ist so kalt, dass es die Haut abschält) für extreme Kälte.

\vspace{3pt}
$\bullet$ Die Temperatur wird in Spanien immer in Celsius angegeben, nie in Fahrenheit.

\columnbreak

$\bullet$ In Südspanien (Andalusien) kann es im Sommer über 40 Grad heiß werden!

\vspace{3pt}
$\bullet$ \es{El barrio} bedeutet "das Viertel" -- jedes spanische Stadtviertel hat seinen eigenen Charakter.

\end{multicols}
\end{grammatikbox}

\end{document}
